\chapter*{Kurzfassung}
\addcontentsline{toc}{chapter}{Kurzfassung}

Obergärige Biere verdanken ihr fruchtiges Aroma vor allem den Estern und höheren
Alkoholen, die Hefen während der Gärung bilden. Diese (erfundene) Arbeit untersucht,
wie stark die Gärtemperatur dieses Aromaprofil prägt. In zwölf Versuchssuden eines
Pale Ale wurde die Temperatur zwischen \SI{16}{\celsius} und \SI{24}{\celsius}
variiert. Die Ergebnisse zeigen einen nahezu linearen Anstieg des Isoamylacetats
-- dem typischen Bananenaroma -- um rund \SI{38}{\percent} je \SI{4}{\kelvin}.
Eine Verkostung mit 24 Personen bestätigt: wärmer vergorene Biere werden als
fruchtiger, aber auch als weniger \enquote{sauber} wahrgenommen.

\chapter*{Abstract}
\addcontentsline{toc}{chapter}{Abstract}

This (fictional) thesis investigates how fermentation temperature shapes the aroma
profile of top-fermented beers. Isoamyl acetate increased by roughly 38\,\% per 4\,K.
